% ===================================================================
%  TABELO N° 6 — L'interna Domo, la Moblaro, la Nutrado, la Lumizado.
%  Feuillets 39 a 47 du fac-simile = folios 37 a 45.
%  Correspondance : folio imprime = numero de feuillet - 2.
%  Le feuillet 48 (folio 46) est un verso blanc, non imprime : la
%  DUESMA SERIO / TABELO N° 7 commence au feuillet 49 (folio 47).
%
%  Les marques « %%K » ne sont pas des commentaires ordinaires : elles
%  portent la CLE D'ALIGNEMENT qui apparie, dans la page de lecture, ce
%  bloc-ci et son homologue francais. La cle se lit
%      t<tableau>-<numero d'alinea numerote>-<rang dans l'alinea>
%  et ne depend d'aucune pagination : les deux editions ne coupent pas
%  leurs pages aux memes endroits, mais elles numerotent les memes
%  alineas — c'est le seul point d'ancrage que les deux livrets
%  partagent, et il est de l'auteur, non du transcripteur.
%
%  Ca tabelo esas dividita en kin chambri (La Dormo-chambro, La
%  Salono, La Manjo-chambro, La Koqueyo, La Balneyo), e l'autoro
%  remetas la numerado di alinei a 1 an la komenco di singla chambro.
%  Onu klavizas ka che videsas, sen korektar la deskresko : outils/
%  ceni.py aplikos pose l'indico di sceno.
% ===================================================================

% --- Feuillet 39 (folio 37, ouverture de tableau) -------------------
\begin{VUpage}[39]{37}
%%K t06-apar-1 apar
\VUsaut{6.0mm}
\VUtitre{15.0pt}{60}{TABELO N\textsuperscript{o} 6}
\VUsaut{1.4mm}
\VUtitre{11.4pt}{0}{L'interna Domo, la Moblaro, la Nutrado,}
\VUsaut{0.4mm}
\VUtitre{11.4pt}{0}{la Lumizado.}
\VUsaut{2.2mm}

%%K t06-apar-2 apar
\VUblancAlinea
La domo esas parfinita. L'onklulo di Ioan\cc
nes esas instalita en sua nova domicilo, di qua\nl
ni konocas l'aranjeso. Nun ni videz quale il\nl
moblizis sua domo. Ni exploros unesme :

%%K t06-tit-1 sub
\VUpk{10.2pt}{40}{La Dormo-chambro.}
\VUsaut{3.0mm}

%%K t06-c1-01-1 p
\VUblancAlinea
1. --- Ni arivas neoportune, nam la cham\cc
bristino esas okupata en netigar la \VUgras{chambro.}

%%K t06-c1-02-1 p
\VUblancAlinea
2. --- Sur la \VUgras{tualet-tablo} \textsuperscript{(1)} esas hike ed ibe\nl
diversa objekti per qui on lavas su, pektas su,\nl
kurtadice on tualetas. Li esas la \VUgras{kuveto} \textsuperscript{(2)} e\nl
l'\VUgras{aquo-krucho} \textsuperscript{(3)}, la \VUgras{kapo-brosilo} \textsuperscript{(4)}, la \VUgras{pek}\cc
\VUgras{tilo} \textsuperscript{(5)}, la \VUgras{sapono} \textsuperscript{(6)} e la \VUgras{sponjo} \textsuperscript{(7)}, fine \VUgras{flako}\cc
\VUgras{neto} \textsuperscript{(8)} por la dentifrico liquida.

%%K t06-c1-03-1 p
\VUblancAlinea
3. --- An la sulo, avan la tualet-tablo, yen\nl
la \VUgras{sitelo} \textsuperscript{(9)} por rekoliar la aqui sordida.

%%K t06-c1-04-1 p
\VUblancAlinea
4. --- Ni vidas en l'\VUgras{avanchambro} \textsuperscript{(10)} kelka\nl
\VUgras{vestohoki} \textsuperscript{(13)} e du \VUgras{parapluvi} \textsuperscript{(11)} en la \VUgras{para}\cc
\VUgras{pluvuyo} \textsuperscript{(12)}.

%%K t06-c1-05-1 p
\VUblancAlinea
5. --- Altralatere di la pordo esas la \VUgras{spegul}\cc
\VUgras{armoro} \textsuperscript{(14)} qua kontenas la \VUgras{linjaro} \textsuperscript{(15)}.

%%K t06-c1-06-1 p
\VUblancAlinea
6. --- Ni profitez ke la \VUgras{chambristino} \textsuperscript{(6)} aran\cc
jas la \VUgras{lito} \textsuperscript{(17)}, por explorar olua diversa parti.\nl
La \VUgras{lito-framo} \textsuperscript{(20)} kontenas bona \VUgras{resort-ma}\cc
\VUgras{traco} \textsuperscript{(21)} sur qua Iustina quik pozos lana \VUgras{ma}\cc
\parplein
\end{VUpage}

% --- Feuillet 40 (folio 38) ------------------------------------------
\begin{VUpage}[40]{38}
%%K t06-c1-06-1 p suite
\VUcontinue
\VUgras{traco} \textsuperscript{(22)}. Pose el prenos de la stuli la du \VUgras{lit}\cc
\VUgras{tuki} \textsuperscript{(23)} e la \VUgras{kovrilo} \textsuperscript{(24)}. Lore el pozos sur la\nl
lit-kapo, la \VUgras{transversa kuseno} \textsuperscript{(25)} la \VUgras{kap-ku}\cc
\VUgras{seno} \textsuperscript{(26)} qui esas kun l'\VUgras{ederduna kovrilo} \textsuperscript{(27)} sur\nl
la \VUgras{tapiso litala} \textsuperscript{(32)}. El uzas plum-fasko por sen\cc
polvigar la \VUgras{kurteni} \textsuperscript{(19)} e la \VUgras{lito-doselo} \textsuperscript{(18)} e yen\nl
parto di la laboro parfinita.

%%K t06-c1-07-1 p
\VUblancAlinea
7. --- Lore el devos ordinar kelke la \VUgras{lito-ta}\cc
\VUgras{bleto} \textsuperscript{(28)}, ritensar la \VUgras{vekigilo} \textsuperscript{(29)}, preparar la\nl
\VUgras{nokto-lampeto} \textsuperscript{(30)}, deprenar la \VUgras{bujio} \textsuperscript{(31)} por\nl
netigar la kandeliero.

%%K t06-tit-2 sub
\VUsaut{5.0mm}
\VUpk{10.2pt}{40}{La laborkorbo e la kamenal acesoraji.}
\VUsaut{3.0mm}

%%K t06-c1-08-1 p
\VUblancAlinea
8. --- La \VUgras{laborkorbo} \textsuperscript{(34)} proxim la \VUgras{tabureto} \textsuperscript{(33)}\nl
anke tre bezonas esar ordinata. El devos faldar\nl
la \VUgras{broduri} \textsuperscript{(35)}, inkluzar en la \VUgras{labor-tablo} \textsuperscript{(38)} la\nl
\VUgras{fingro-chapo}, la \VUgras{filo} \textsuperscript{(36)}, l'\VUgras{aguli} \textsuperscript{(37)} e la \VUgras{cizo} \textsuperscript{(39)}\nl
e klozar per la \VUgras{klefo} \textsuperscript{(40)} la kovrilo di la tablo.

%%K t06-c1-09-1 p
\VUblancAlinea
9. --- Sur la \VUgras{unpeda tablo} \textsuperscript{(41)} dil mezo, Ius\cc
tina rinovigos l'aquo di la \VUgras{karafo} \textsuperscript{(42)}. El pozos\nl
\VUgras{sukro} en la \VUgras{sukruyo} \textsuperscript{(43)}, netigos la \VUgras{sukro-pre}\cc
\VUgras{nilo} \textsuperscript{(44)} e deprenos la \VUgras{vesto-brosilo} \textsuperscript{(46)}.

%%K t06-c1-10-1 p
\VUblancAlinea
10. --- La dextra latero di la chambro postu\cc
los de elu min multa laboro, nam la \VUgras{longa stu}\cc
\VUgras{lego} \textsuperscript{(47)} e lua \VUgras{kuseno} \textsuperscript{(48)} esas en sua justa plaso,\nl
e, pro ke ne esas kolda vetero, nulo esas dipla\cc
sita inter la acesoraji (implementi) di la \VUgras{ka}\cc
\VUgras{meno} \textsuperscript{(49)}. \VUgras{Shovelo} \textsuperscript{(50)}, \VUgras{pincegi} \textsuperscript{(51)}, \VUgras{morilii} \textsuperscript{(55)},\nl
\VUgras{cindro-barilo} \textsuperscript{(56)} esas neta; la \VUgras{fairo-skreno} \textsuperscript{(54)}\nl
ne esis diplasita e la \VUgras{ligno-kesto} \textsuperscript{(52)} esas bone\nl
provizita per \VUgras{lenii} \textsuperscript{(53)}.
\end{VUpage}

% --- Feuillet 41 (folio 39) -------------------------------------------
%  Cette page porte DEUX notes (asterisque page-local, remis a 1 pour
%  chaque appel) : l'une glose « babo », l'autre « lambrequino ».
\begin{VUpage}[41]{39}
\VUnotes{143.2mm}{%
%%K t06-noto-1 noto
(*) Babo = infanteto, la E. baby, F. bébé, I. bambino.
\par
%%K t06-noto-2 noto
(*) Lambrequino = F. lambrequin, S. lambrequines, Port.\nl
lambrequins, mem D. E. lambrequins, do D. E. F. P. S.%
}

%%K t06-c1-11-1 p
\VUblancAlinea
11. --- Tamen el devos senpolvigar la \VUgras{pen}\cc
\VUgras{dul-horlojo} \textsuperscript{(57)}, la \VUgras{kandelabri} \textsuperscript{(58)} e la \VUgras{posh}\cc
\VUgras{kozuyi} \textsuperscript{(59)}, fine vishar la \VUgras{spegulo} \textsuperscript{(60)}.

%%K t06-c1-12-1 p
\VUblancAlinea
12. --- Pri la \VUgras{bersilo} \textsuperscript{(61)} dil infanteto (dil\nl
babo (*)), ol ja esas pronta.

%%K t06-tit-3 sub
\VUsaut{5.0mm}
\VUpk{10.2pt}{40}{La Salono.}
\VUsaut{3.0mm}

%%K t06-c2-01-1 p
\VUblancAlinea
1. --- Nun yen la salono. Lu esas tre bela\nl
\VUgras{chambro.} La kameno, qua situesas an la dex\cc
tra angulo esas ornita per delikata \VUgras{statueto} \textsuperscript{(1)},\nl
per du bela \VUgras{vazi} \textsuperscript{(3)} e drapirita per velura \VUgras{lam}\cc
\VUgras{brequino} \textsuperscript{(2)} (*).

%%K t06-c2-02-1 p
\VUblancAlinea
2. --- Por impedar ke la cintili di la herdo\nl
bruleskez la tapiso, on pozis avan la herdo ku\cc
pra \VUgras{paracintilo} \textsuperscript{(4)}.

%%K t06-c2-03-1 p
\VUblancAlinea
3. --- Proxime, eleganta \VUgras{skribo-tablo} \textsuperscript{(5)} siori\cc
nala esas apertita. Sur olu la \VUgras{domo-mastri}\cc
\VUgras{no} \textsuperscript{(23)} skribas sua letri.

%%K t06-c2-04-1 p
\VUblancAlinea
4. --- La fenestro esas garnisita per \VUgras{vitral}\nl
\VUgras{kurteni} \textsuperscript{(6)} kun \VUgras{ligili} \textsuperscript{(8)} akrochita al \VUgras{pateri} \textsuperscript{(7)}\nl
e per stofa kurtenegi \VUgras{supre drapirita} \textsuperscript{(9)}.

%%K t06-c2-05-1 p
\VUblancAlinea
5. --- En l'embrazuro di la fenestro, \VUgras{po}\cc
\VUgras{tuyo} \textsuperscript{(11)} kontenas verda \VUgras{planto} \textsuperscript{(10)}, splendida\nl
palmiero di qua la folii su extensas preske til la\nl
\VUgras{petrol-lampo} \textsuperscript{(12)}.

%%K t06-c2-06-1 p
\VUblancAlinea
6. --- La \VUgras{lum-shirmilo} \textsuperscript{(13)} di la lampo\nl
esas aranjita da Maria, la charmanta yuni\cc
\parplein
\end{VUpage}

% --- Feuillet 42 (folio 40) -------------------------------------------
\begin{VUpage}[42]{40}
%%K t06-c2-06-1 p suite
\VUcontinue
no \textsuperscript{(14)} qua esas an la piano \textsuperscript{(15)}. Tre rekte si\cc
danta sur la \VUgras{tabureto} \textsuperscript{(16)}, el studias muzikajo.\nl
El esas tre habila e sorgema, quale pruvas\nl
elua \VUgras{muzik-armoro} \textsuperscript{(17)} qua esas perfekte ordi\cc
nita.

%%K t06-tit-4 sub
\VUsaut{5.0mm}
\VUpk{10.2pt}{40}{La Bubacho.}
\VUsaut{3.0mm}

%%K t06-c2-07-1 p
\VUblancAlinea
7. --- Elua fratulo, Paulus, serchas \VUgras{libro} \textsuperscript{(19)}\nl
en la \VUgras{biblioteko} \textsuperscript{(18)}. Il esas \VUgras{yuno} \textsuperscript{(20)}, movema\nl
e netolerebla. Akrochante su al biblioteko por\nl
atingar la libro qua esas tro supra, il riskas\nl
faligar la \VUgras{busto} \textsuperscript{(21)} qua esas supere.

%%K t06-c2-08-1 p
\VUblancAlinea
8. --- Il profitas, por facar ta stultajo, ke\nl
sua matro esas okupata per konversar kun\nl
amikino, \VUgras{damo} \textsuperscript{(22)} qua venis pasar kelka dii\nl
en lua hemo. La matro sidas sur \VUgras{kanapeo} \textsuperscript{(24)}\nl
proxim la \VUgras{stulego} \textsuperscript{(25)}, ed elua amikino esas sur\nl
la \VUgras{sidilo} \textsuperscript{(27)} di qua on vidas nur la \VUgras{dors-apo}\cc
\VUgras{gilo} \textsuperscript{(26)}.

%%K t06-c2-09-1 p
\VUblancAlinea
9. --- La granda \VUgras{kadro} \textsuperscript{(30)} pendanta an la\nl
muro kontenas la \VUgras{portreto} \textsuperscript{(29)} di parentino.\nl
En l'\VUgras{albumo} \textsuperscript{(32)} pozita sur la tablo esas asem\cc
blita la \VUgras{fotografuri} \textsuperscript{(33)} di la cetera familiani.\nl
La pueri jus foliumis olu, nam la \VUgras{stuli} \textsuperscript{(31)} esas\nl
ankore grupigita cirkum la tablo.

%%K t06-c2-10-1 p
\VUblancAlinea
10. --- La \VUgras{Chiniana vazo} \textsuperscript{(34)} porcelana, qua\nl
esas proxim la \VUgras{paper-kultelo} \textsuperscript{(35)} esis donacita\nl
a Maria por elua festo, e la \VUgras{paravento} \textsuperscript{(36)} qua\nl
garnisas l'angulo di la chambro esas adportita\nl
de Chinia da elua onklulo.

%%K t06-c2-11-1 p
\VUblancAlinea
11. --- Gayigata dal belega \VUgras{pikturi} \textsuperscript{(37)} akro\cc
chita an la \VUgras{parieto} \textsuperscript{(42)}, lumizata dal \VUgras{plafon}\cc
\VUgras{lustro} \textsuperscript{(38)}, di qua l'\VUgras{elektral-lampi} \textsuperscript{(39)} vespere
\parplein
\end{VUpage}

% --- Feuillet 43 (folio 41) -------------------------------------------
\begin{VUpage}[43]{41}
%%K t06-c2-11-1 p suite
\VUcontinue
brilas joyoze, ta salono aspektas tre agreabla.\nl
La mola \VUgras{tapiso} \textsuperscript{(40)} extensita sur la parqueto, e\nl
la tante faciluza \VUgras{timbro elektrala} \textsuperscript{(41)}, parfinas\nl
olua komforto.

%%K t06-tit-5 sub
\VUsaut{3.6mm}
\VUpk{10.2pt}{40}{La Manjo-chambro.}
\VUsaut{3.0mm}

%%K t06-c3-01-1 p
\VUblancAlinea
1. --- Sonis la dineo-kloko. La tota familio,\nl
ecepte Maria, esas asemblita cirkum la tablo.\nl
La manjo-chambro esas agreable varmigita dal\nl
bonega \VUgras{furnelo} \textsuperscript{(1)} fayenca, kun tenua \VUgras{tubo} \textsuperscript{(2)}.

%%K t06-c3-02-1 p
\VUblancAlinea
2. --- Ta chambro kontenas ne multa mobli.\nl
Alonge la muro pendas, super la \VUgras{skabelo} \textsuperscript{(4)}.\nl
\VUgras{fonteneto} \textsuperscript{(3)} orniva. Tre proxime, esas la \VUgras{kre}\cc
\VUgras{denco} \textsuperscript{(5)} sur qua on pozas la kolda manjebli,\nl
la deser-dishi, e la diversa tabl-implementi\nl
quin on ne quik bezonas.

%%K t06-c3-03-1 p
\VUblancAlinea
3. --- Tale ni vidas sur la supera tabulo,
\parplein
\end{VUpage}

% --- Feuillet 44 (folio 42) -------------------------------------------
\begin{VUpage}[44]{42}
%%K t06-c3-03-1 p suite
\VUcontinue
\VUgras{hanyuno rostita} \textsuperscript{(7)}, \VUgras{fisho} \textsuperscript{(8)} e \VUgras{kupo} \textsuperscript{(10)} kun\nl
\VUgras{frukti} \textsuperscript{(9)}. Sube yen la \VUgras{fromajo} \textsuperscript{(11)}, la \VUgras{pano} \textsuperscript{(12)},\nl
la \VUgras{flakonuyo} \textsuperscript{(13)} qua kontenas omno bezonata\nl
por kondimentizar la salado : l'oleo, la vina\cc
gro, la \VUgras{salo} \textsuperscript{(14)} e la \VUgras{pipro} \textsuperscript{(15)}. Fine, sur l'infra\nl
tabulo esas \VUgras{kuko} \textsuperscript{(17)} proxim la \VUgras{mustarduyo} \textsuperscript{(16)}.

%%K t06-c3-04-1 p
\VUblancAlinea
4. --- La manjo-chambrala horlojo, quan on\nl
nomizas \VUgras{mur-horlojo} \textsuperscript{(20)}, havas tre aparta\nl
formo. Ol esas enkastrita en ligna kadro qua\nl
pluse kontenas \VUgras{barometro} \textsuperscript{(19)} e \VUgras{termometro} \textsuperscript{(18)}.

%%K t06-tit-6 sub
\VUsaut{4.6mm}
\VUpk{10.2pt}{40}{La Dineo-servo.}
\VUsaut{3.0mm}

%%K t06-c3-05-1 p
\VUblancAlinea
5. --- An omna muri di la chambro esas\nl
aplikita \VUgras{ligno-paneli} \textsuperscript{(21)} ek vaxizita querko.\nl
Stofa \VUgras{pordo-tapeto} \textsuperscript{(22)} sate klozas la pordo\nl
qua donas eniro por la verando. Ibe la fami\cc
lio iros drinkar la kafeo pos la dineo. Ja la\nl
\VUgras{tasi} \textsuperscript{(24)} e la \VUgras{subtasi} \textsuperscript{(25)} esas pronta sur la\nl
\VUgras{plad-moblo} \textsuperscript{(23)}.

%%K t06-c3-06-1 p
\VUblancAlinea
6. --- Super la tablo esas fixigita al plafono\nl
\VUgras{pendo-lumizilo} \textsuperscript{(26)} di qua la tre brilanta lampo\nl
vespere lumizas la tota manjo-chambro.

%%K t06-c3-07-1 p
\VUblancAlinea
7. --- Nun ni explorez la \VUgras{manjo-tablo} \textsuperscript{(27)}\nl
ipsa. Kande eniris la familio, la manjilaro esis\nl
pronta. Sur la \VUgras{tablo-tuko} \textsuperscript{(28)} esis pozita, ye\nl
la plaso di singla kunmanjonto, \VUgras{plado} \textsuperscript{(33)},\nl
\VUgras{bok-tuko} \textsuperscript{(29)}, \VUgras{kuliero} \textsuperscript{(34)}, \VUgras{forketo} \textsuperscript{(35)}, \VUgras{kulte}\cc
\VUgras{lo} \textsuperscript{(36)} e \VUgras{glaso} \textsuperscript{(32)}. Esis anke \VUgras{boteli} \textsuperscript{(30)} por vino\nl
e \VUgras{karafo} \textsuperscript{(31)} por aquo.

%%K t06-c3-08-1 p
\VUblancAlinea
8. --- La \VUgras{servisto} \textsuperscript{(39)} unesme adportis la \VUgras{su}\cc
\VUgras{puyo} \textsuperscript{(37)}, en qua ankore fumas kelka supo.\nl
Nun il servas l'unesma \VUgras{disho} \textsuperscript{(38)}, karno-disho.\nl
Pose il prizentos olti quin ni ja nomis, fine,\nl
pos la \VUgras{deser-disho}, il profitos ke sua mastri\nl
enirabos la verando, por deprenar la tabl-uten\cc
sili e riordinar la manjo-chambro.

%%K t06-c3-08-2 p
\VUblancAlinea
\textit{Noto.} --- \textit{La komunuza vorti relatanta la nu}\cc
\textit{trivaro, hiké indikesas tre rezume. La listo}\nl
\textit{esos kompletigata sur la du tabeli « La Hote}\cc
\textit{lo » n\textsuperscript{o} 13) e « La Merkato » (n\textsuperscript{o} 15).}

%%K t06-tit-7 sub
\VUsaut{4.6mm}
\VUpk{10.2pt}{40}{La Koqueyo.}
\VUsaut{3.0mm}

%%K t06-c4-01-1 p
\VUblancAlinea
1. --- En la koqueyo, quan ni regardeskas\nl
tra la pordo mi-apertita, ni vidas pitoreska\nl
desordino. Avan la \VUgras{herdo} \textsuperscript{(1)} en qua \VUgras{fairo} \textsuperscript{(3)}\nl
krepitas, esas instalita \VUgras{spiso-turnilo} \textsuperscript{(5)}. Bela\nl
\VUgras{rostajo} \textsuperscript{(7)} esas \VUgras{spisagita} \textsuperscript{(6)} e parkoquas.
\end{VUpage}

% --- Feuillet 45 (folio 43) -------------------------------------------
\begin{VUpage}[45]{43}
%%K t06-c4-02-1 p
\VUblancAlinea
2. --- La \VUgras{suflilo} \textsuperscript{(8)} e la \VUgras{karbon-shovelo} \textsuperscript{(9)},\nl
quin on jus uzis, restis an la sulo same kam\nl
la \VUgras{lenii} \textsuperscript{(10)}.

%%K t06-c4-03-1 p
\VUblancAlinea
3. --- La suprajo di la kameno esas ordinita;\nl
la \VUgras{lanterno} \textsuperscript{(11)}, l'\VUgras{alumetuyo} \textsuperscript{(13)}, en qua esas\nl
l'\VUgras{alumeti} \textsuperscript{(12)}, la \VUgras{kandeliero} \textsuperscript{(14)} e la \VUgras{bujiero} \textsuperscript{(15)}\nl
kun \VUgras{bujio} \textsuperscript{(16)} esas en sua plaso. Ma la granda\nl
\VUgras{karbon-sitelo} \textsuperscript{(18)}, plena de \VUgras{terkarbono} \textsuperscript{(17)} quan\nl
la \VUgras{koquistino} \textsuperscript{(23)} jus plenigis en la kelero, res\cc
tis meze di la koqueyo.

%%K t06-c4-04-1 p
\VUblancAlinea
4. --- Advere, la kompatinda muliero devas\nl
hastar por ke la dejuno esez pronta justatempe.\nl
Nun el esas proxim la \VUgras{koquo-furnelo} \textsuperscript{(19)} e\nl
quirlas \VUgras{sauco} \textsuperscript{(24)} quan oportas ne brular troe.

%%K t06-c4-05-1 p
\VUblancAlinea
5. --- En la \VUgras{forno} \textsuperscript{(20)} bakas kuko, quan el\nl
preparis por la repasteto dil pueri. Super la\nl
koquo-furnelo pendas la \VUgras{pot-kuliero} \textsuperscript{(21)} e la\nl
\VUgras{spum-kuliero} \textsuperscript{(22)}.

%%K t06-tit-8 sub
\VUsaut{4.0mm}
\VUpk{10.2pt}{40}{La Vazaro.}
\VUsaut{3.0mm}

%%K t06-c4-06-1 p
\VUblancAlinea
6. --- La \VUgras{koquilaro} \textsuperscript{(25)}, akrochita funde di la\nl
chambro, esas tre neta. La bel kupra \VUgras{kas}\cc
\VUgras{roli} \textsuperscript{(26)}, la \VUgras{lakto-poto} \textsuperscript{(27)}, la \VUgras{kribleto} \textsuperscript{(28)} e la\nl
\VUgras{raspilo} \textsuperscript{(29)} esis sorgoze netigita.

%%K t06-c4-07-1 p
\VUblancAlinea
7. --- La \VUgras{sal-buxo} \textsuperscript{(30)} esas pozita sur la plad-\nl
armoro proxim la \VUgras{te-krucho} \textsuperscript{(31)}, e l'\VUgras{oleo-bo}\cc
\VUgras{telego} \textsuperscript{(32)}. La \VUgras{tirkesto} \textsuperscript{(33)} kredeble kontenas la\nl
manjilari e la kulteli di koqueyo.

%%K t06-c4-08-1 p
\VUblancAlinea
8. --- La \VUgras{balayilo} \textsuperscript{(34)} e la \VUgras{plum-fasko} \textsuperscript{(35)} esas\nl
en angulo. La koquistino jus uzis la \VUgras{ligno}\cc
\VUgras{bloko} \textsuperscript{(36)}, el lasis ol proxim la fenestro.
\end{VUpage}

% --- Feuillet 46 (folio 44) -------------------------------------------
\begin{VUpage}[46]{44}
%%K t06-c4-09-1 p
\VUblancAlinea
9. --- \VUgras{Manuvishilo} \textsuperscript{(37)} pendas an la muro,\nl
proxim la \VUgras{funelo} \textsuperscript{(38)}. En ta parto di la ko\cc
queyo esas forpozita la \VUgras{rostilo} \textsuperscript{(39)}, la \VUgras{ha}\cc
\VUgras{chilo} \textsuperscript{(40)} e la \VUgras{hachoplanko} \textsuperscript{(41)}. Pose, super la\nl
hachilo, sur \VUgras{tabulo} \textsuperscript{(42)}, esas \VUgras{poti} \textsuperscript{(43)}, la \VUgras{bolio}\cc
\VUgras{poto} \textsuperscript{(44)} kun sua \VUgras{kovrilo} \textsuperscript{(45)} e la \VUgras{marmito} \textsuperscript{(46)}.\nl
La \VUgras{padelo} \textsuperscript{(47)} pendas proxime.

%%K t06-c4-10-1 p
\VUblancAlinea
10. --- Nur la kupra \VUgras{robineto} \textsuperscript{(48)} projektas\nl
brilo en ica angulo. L'\VUgras{eviero} \textsuperscript{(49)} sur qua esas\nl
pozita la dika \VUgras{kruchego} \textsuperscript{(50)}, esas kredeble tre\nl
komoda kun sua mikra armoro, en qua on\nl
konservas la \VUgras{vinagro-barelo} \textsuperscript{(51)} provizita per\nl
sua \VUgras{robineto} \textsuperscript{(52)}, la \VUgras{rezidu-buxo} \textsuperscript{(53)} et la \VUgras{des}\cc
\VUgras{neta pladi} \textsuperscript{(54)}.

%%K t06-tit-9 sub
\VUsaut{4.0mm}
\VUpk{10.2pt}{40}{Altra objekti pozita en la koqueyo.}
\VUsaut{3.0mm}

%%K t06-c4-11-1 p
\VUblancAlinea
11. --- La \VUgras{kuvego} \textsuperscript{(55)} en qua on lavas la \VUgras{le}\cc
\VUgras{gumi} \textsuperscript{(58)} et la gisfera \VUgras{kaldrono} \textsuperscript{(56)} pozita sur la\nl
\VUgras{karel-sulo} \textsuperscript{(57)} kredeble havas anke sua plaso\nl
en ta armoro.

%%K t06-c4-12-1 p
\VUblancAlinea
12. --- La koquistino certe ne indijas furneli,\nl
nam yen pluse, sur l'angulo di la tablo, \VUgras{gas}\cc
\VUgras{furnelo} \textsuperscript{(59)} sur qua varmeskas aquo en mikra\nl
kasrolo. La \VUgras{vaporo} \textsuperscript{(60)} qua eskapas ek olu indi\cc
kas a ni, ke kredeble ol bolias. Versimile ca\nl
aquo uzesos por la kafeo, nam yen sur l'altra\nl
extremajo dil tablo, la \VUgras{kafe-krucho} \textsuperscript{(64)} e la\nl
\VUgras{kafe-muelilo} \textsuperscript{(63)}.

%%K t06-c4-13-1 p
\VUblancAlinea
13. --- La furnelo esas konektata kun la \VUgras{gas}\cc
\VUgras{brulilo} \textsuperscript{(62)} per longa \VUgras{kauchuka tubeto} \textsuperscript{(61)}.

%%K t06-c4-14-1 p
\VUblancAlinea
14. --- La \VUgras{jornala servistino} \textsuperscript{(66)} qua venas\nl
helpar la koquistino, preparas la legumi por
\parplein
\end{VUpage}

% --- Feuillet 47 (folio 45) -------------------------------------------
\begin{VUpage}[47]{45}
%%K t06-c4-14-1 p suite
\VUcontinue
la dineo. Kande oli esabos emundata e lavata,\nl
el pozos li en la \VUgras{baseno} \textsuperscript{(65)} vartante la horo por\nl
koquar li.

%%K t06-tit-10 sub
\VUsaut{4.0mm}
{\centering\fontsize{10.2pt}{10.2pt}\selectfont\textls[40]{\scshape La Balneyo.} \textit{(Balno-chambro.)}\par}
\VUsaut{3.0mm}

%%K t06-c5-01-1 p
\VUblancAlinea
1. --- A ni restas facenda nur un nediskre\cc
tajo por konoceskar la precipua chambri dil\nl
apartamento : vidar l'instaluro di la balneyo.\nl
To ne postulos longa tempo, nam ol ne esas\nl
granda, e ni rapide savos ke la \VUgras{balno-kuvo} \textsuperscript{(1)}\nl
havas bona \VUgras{balno-varmigilo} \textsuperscript{(2)}, ke la \VUgras{sapon}\cc
\VUgras{konko} \textsuperscript{(3)} esas atingebla da la manuo di bal\cc
nanto e ke la \VUgras{tuk-portilo} \textsuperscript{(5)} certe igos facile\nl
sikigar la \VUgras{tualet-tuki} \textsuperscript{(4)}.

%%K t06-c5-02-1 p
\VUblancAlinea
2. --- Ta chambro havas sua parieti kovrita\nl
per \VUgras{fayencaji} \textsuperscript{(6)}, qui posibligis instalar \VUgras{dush}\cc
\VUgras{aparato} \textsuperscript{(7)} sen irge timar ke l'aquo, qua retro\cc
spricas dum lua uzeso, domajos la muri. Zinka\nl
\VUgras{baseneto} \textsuperscript{(8)} qua servas por la ped-balni, kom\cc
pletigas la moblaro.

\VUsaut{6.0mm}
\VUfilet{22mm}
\end{VUpage}

% --- Feuillet 48 (folio 46) : verso blanc, non imprime. -------------
